\section{Results}
\label{sec:results}

The results are reported in the order of the analyses described in
Section~\ref{sec:methods}. First, the discrimination of the primary ensemble and its
transfer to unseen sites are presented, together with a placement against published
ABIDE-I work. Then, the six contributions are reported in turn. In order, these concern
cognitive ability, functional-system differences under motion matching, feature-selection
scope, classifier family, the scope of the symptom-severity associations, and the covariance
shrinkage applied by the windowed description. Finally, the secondary results are summarised,
and their null findings are stated as limits. Throughout, AUROC and accuracy are labelled as
participant-pooled or as fold means (Section~\ref{subsubsec:performance-measures}), because the two summaries
\reviewhl{differ}.

\subsection{Classification performance and transfer to unseen sites}
\label{subsec:results-performance}

The highest AUROC of the connectivity-only configurations is obtained with the primary
ensemble under both validation protocols (Table~\ref{tab:configurations}). Specifically, a
fold-mean AUROC of 0.775 is obtained under pooled ten-fold validation, and a fold-mean
AUROC of 0.773 is obtained under leave-one-site-out validation. However, the gain over the
best single representation, the windowed tangent mean, is small ($+0.009$ in
participant-pooled AUROC).
Moreover, this gain is not distinguished from chance (unadjusted DeLong $p=0.066$), so it
is not presented as an established improvement. The rank-mean ensemble
(Section~\ref{subsubsec:rank-fusion}) is, in turn, reported for comparison only. Because its score is
not a probability, it has no 0.5 decision threshold, and no accuracy is reported for it.

\begin{table}[pos=htb]
\centering
\caption{Discrimination of the ensemble and of its members in the primary cohort
(871 participants). AUROC and accuracy are fold means, and SD is the standard deviation
across outer folds. The rank-mean ensemble has no probability scale, so no accuracy is
given for it.}
\label{tab:configurations}
\small
\begin{tabular}{@{}lcccccc@{}}
\toprule
 & \multicolumn{3}{c}{Pooled ten-fold} & \multicolumn{3}{c}{Leave-one-site-out} \\
\cmidrule(lr){2-4}\cmidrule(l){5-7}
Configuration & AUROC & SD & Accuracy & AUROC & SD & Accuracy \\
\midrule
Ensemble, probability mean (adopted) & 0.775 & 0.040 & 0.707 & 0.773 & 0.087 & 0.684 \\
Windowed tangent, mean & 0.767 & 0.044 & 0.695 & 0.771 & 0.082 & 0.670 \\
Ensemble, rank mean & 0.756 & 0.038 & -- & 0.750 & 0.091 & -- \\
Static tangent & 0.756 & 0.048 & 0.692 & 0.757 & 0.104 & 0.643 \\
Correlation & 0.750 & 0.040 & 0.683 & 0.743 & 0.083 & 0.655 \\
\bottomrule
\end{tabular}
\end{table}

As shown by participant-pooled estimates with intervals, discrimination is maintained when
a whole site is withheld from training (Table~\ref{tab:performance};
Figure~\ref{fig:performance}a). Specifically, an AUROC of 0.776 (95\% DeLong interval
0.745--0.807) is obtained under pooled ten-fold validation. Similarly, an AUROC of 0.774
(0.744--0.805) is obtained under leave-one-site-out validation. In addition, sensitivity
and specificity are nearly equal under pooled ten-fold validation (0.705 and 0.709), so the
ensemble does not favour one class under this protocol. With the site withheld, in
contrast, sensitivity is higher than specificity (0.742 against 0.684). As a result, more of
the errors under this protocol fall on non-autistic participants. These intervals, however,
condition on the fitted model and treat
dependent out-of-fold predictions as independent (Section~\ref{subsubsec:interval-construction}). For this
reason, they are the narrowest defensible intervals rather than a full account of
uncertainty.

Intervals that carry the site clustering are wider, and they are reported alongside the
parametric ones (Section~\ref{subsubsec:paired-comparisons}). Under the site-clustered
bootstrap, the participant-pooled AUROC of 0.776 carries a 95\% interval of 0.733 to 0.807,
against 0.745 to 0.807 from the DeLong interval. Similarly, accuracy carries an interval of
0.670 to 0.737, against 0.676 to 0.736 from the Wilson interval. With a site withheld, the
AUROC interval widens from 0.744--0.805 to 0.726--0.808, and the accuracy interval from
0.680--0.740 to 0.661--0.748. The lower bound moves further than the upper bound in every
case, so the parametric intervals are anti-conservative mainly at the lower end. None of the
intervals includes 0.5, so the discrimination itself is not placed in question by the wider
estimates.

The analysed cohorts are reconciled against the released records. Of the 1,112 ABIDE-I
phenotypic records, 1,035 have a preprocessed derivative, and 871 satisfy the three-rater
primary rule and have usable regional time series, comprising 403 autistic and 468
non-autistic participants.

\begin{table}[t]
\centering
\caption{Participant-pooled performance of the primary ensemble in the primary cohort
(871 participants), with 95\% intervals. AUROC carries DeLong intervals; accuracy,
sensitivity and specificity carry Wilson intervals.}
\label{tab:performance}
\small
\begin{tabular}{@{}lcccc@{}}
\toprule
Validation & AUROC & Accuracy & Sensitivity & Specificity \\
\midrule
Pooled ten-fold & 0.776 (0.745--0.807) & 0.707 (0.676--0.736) & 0.705 (0.658--0.747) & 0.709 (0.667--0.749) \\
Leave-one-site-out & 0.774 (0.744--0.805) & 0.711 (0.680--0.740) & 0.742 (0.697--0.782) & 0.684 (0.640--0.724) \\
\bottomrule
\end{tabular}
\end{table}

\begin{figure}[pos=tbp]
\centering
\includegraphics[width=\linewidth]{fig_performance}
\caption{Discrimination, per-site transfer, calibration and abstention of the primary
ensemble. (a) Discrimination with 95\% intervals under both validation protocols. (b) AUROC
at every site when that site is withheld from training, with marker area proportional to
the number of participants; the random-effects mean and the prediction interval for a site
not yet seen are shown. (c) Predicted probability against observed frequency. (d) For the
cohort-wide and per-site abstention rules under leave-one-site-out validation, the
percentage of participants answered when the threshold is set on the whole sample and the
percentage of held-out half-splits in which the error target is met, against the nominal
90\%.}
\label{fig:performance}
\end{figure}

Variation between sites is summarised by a random-effects model \reviewhl{as a descriptive sensitivity analysis}. When the 20 per-site
estimates are combined, therefore, a random-effects mean AUROC of 0.803 is obtained
(Figure~\ref{fig:performance}b). In addition, between-site heterogeneity is not
distinguishable from chance ($Q=23.5$ on 19 degrees of freedom, $p=0.22$; $I^2=19\%$). As a
result, the \reviewhl{descriptive 95\% interval} of Equation~\eqref{eq:prediction-interval}
extends from 0.725 to 0.882. However, leave-one-site-out estimates share training data, so
the heterogeneity is, if anything, understated. \reviewhl{This interval is not interpreted as a validated prediction interval for an independent site.}

Site-level AUROC nevertheless varies widely, from 0.589 at SBL to 0.912 at UM\_2. In
particular, no site's point estimate falls to chance, but 2 of 20 sites lie below 0.65
(SBL and MAX\_MUN). Moreover, the 95\% intervals of 5 sites (CALTECH, LEUVEN\_1, MAX\_MUN,
OHSU and SBL) include 0.5. This spread, however, is not taken as evidence of site-dependent
failure, because the median site contributes about thirty participants. For this reason,
wide per-site intervals are expected even under identical performance at every site.

Phenotypic variables carry some diagnostic information, but they add little to the
connectivity features. Specifically, an AUROC of about 0.62 is reached by the phenotype
model (Section~\ref{par:phenotype-model}) on its own. In addition, participant-pooled AUROC is raised
by about 0.004 when this model is added to the ensemble (unadjusted $p=0.041$; $q=0.11$,
for reference only). Because this gain is small, the phenotypic member is not adopted.

The placement against published ABIDE-I work is given in Table~\ref{tab:published-comparison}, in
which sample size and testing procedure are listed because both change what a value means.
For this comparison, the wider quality-control cohort is used,
because its quality-control definition admits 1,035 participants, the same eligible count
as in the published studies. However, three of these participants have no usable time
series in the release, so 1,032 are analysed. In this cohort, a participant-pooled accuracy
of 0.706 (95\% Wilson interval 0.678--0.733) is obtained under pooled ten-fold validation. In comparison, accuracies of 0.700, 0.703 and 0.708 are
reported by \citet{Heinsfeldetal2018}, \citet{Eslamietal2019} and
\citet{AlmuqhimSaeed2021}, and these values are not recomputed here.

\begin{table}[t]
\centering
\caption{Placement against published work on ABIDE-I resting-state functional
connectivity. Values for the other studies are as reported by their authors; n.r., not
reported. For the ensemble rows, AUROC is participant-pooled and accuracy is the fold mean,
which agrees with the participant-pooled value to three decimals. For the tangent-model
row, AUROC and accuracy are fold means. The unseen-site column gives the
leave-one-site-out AUROC, with its difference from the pooled ten-fold value in brackets.}
\label{tab:published-comparison}
\small
\setlength{\tabcolsep}{4pt}
\begin{tabular}{@{}>{\raggedright\arraybackslash}p{0.24\linewidth}cccccc@{}}
\toprule
Study & Accuracy & AUROC & Sample & Procedure & AUROC, unseen site & Calibration error \\
\midrule
\citet{Heinsfeldetal2018} & 0.700 & n.r. & 1,035 & 10-fold & n.r. & n.r. \\
\citet{Eslamietal2019} & 0.703 & n.r. & 1,035 & 10-fold & n.r. & n.r. \\
\citet{AlmuqhimSaeed2021} & 0.708 & n.r. & 1,035 & 10-fold & n.r. & n.r. \\
This work, ensemble & 0.707 & 0.776 & 871 & nested 10-fold & 0.774 ($-0.002$) & 0.046 \\
This work, ensemble, \reviewhl{wider QC cohort} & 0.706 & 0.789 & 1,032 & nested 10-fold & 0.783 ($-0.006$) & 0.023 \\
This work, tangent model, \reviewhl{wider QC cohort} & 0.696 & 0.766 & 1,032 & nested 10-fold & -- & -- \\
\bottomrule
\end{tabular}
\end{table}

The accuracy comparison in Table~\ref{tab:published-comparison} does not separate the methods,
because none of the comparison studies reports a confidence interval. In addition, the
intervals implied by these sample sizes are roughly $\pm0.03$, whereas the three published
values span 0.008. For this reason, two further quantities that those studies do not
report are listed, and they describe the present work without being compared with the
published values. Specifically, the AUROC on an unseen site is 0.774 against 0.776 within
the cohort, and the expected calibration error is 0.046 (Section~\ref{subsec:results-secondary}). In
addition, accuracy at the published level is reached under a procedure in which no step is
fitted before the data are split. This property matters because a step fitted before the
split \reviewhl{inflates} the estimate (Section~\ref{subsec:results-selection-scope}).

The cohort definition itself changes the score, so both definitions are reported for a
like-for-like comparison. Specifically, a fold-mean AUROC of 0.756 is reached by the static
tangent model in the primary cohort of 871 participants, and 0.766 in the wider cohort of
1,032 participants. However, the scans excluded by the primary definition are mostly those with
head motion or failed processing, and they are not evenly split between the groups. Part of
this gain \reviewhl{could}, therefore, reflect data quality rather than differences in brain
connectivity.

\subsection{Cognitive ability and detection}
\label{subsec:results-cognitive-ability}

\reviewhl{Discrimination is lower in higher-IQ autistic subgroups in the unmatched band analysis}, and
this pattern underlies the first contribution of the study. Because the subgroup analyses
are exploratory (Sections~\ref{par:cognitive-sex-subgroups} and~\ref{subsubsec:analysis-status}), the pattern is reported as
a description rather than as a confirmed effect. To separate it from changes in the
comparison group, the comparison group is held fixed at all 468 non-autistic participants.
Only the autistic participants, in turn, are divided into full-scale IQ bands. As a result, every band is
scored against the same reference, so differences between bands cannot arise from changes
in the comparison group. Under this design, AUROC falls monotonically from 0.873 below
full-scale IQ 85 to 0.719 at 115 or above under pooled ten-fold validation
(Table~\ref{tab:iq-bands}; Figure~\ref{fig:group-differences}d). Similarly, AUROC falls from
0.871 to 0.709 when the site is withheld. However, the autistic bands still differ from the
reference group in full-scale IQ, and this separation is addressed next.

\begin{table}[t]
\centering
\caption{AUROC by full-scale IQ (FIQ) band of the autistic participants, each band scored
against the same comparison group of all 468 non-autistic participants.}
\label{tab:iq-bands}
\small
\begin{tabular}{@{}lccc@{}}
\toprule
Autistic participants & $n$ & Pooled ten-fold & Leave-one-site-out \\
\midrule
FIQ below 85 & 40 & 0.873 & 0.871 \\
FIQ 85--99 & 97 & 0.815 & 0.816 \\
FIQ 100--114 & 126 & 0.781 & 0.788 \\
FIQ 115 or above & 113 & 0.719 & 0.709 \\
\bottomrule
\end{tabular}
\end{table}

\begin{figure}[pos=tbp]
\centering
\includegraphics[width=\linewidth]{fig_groupdiff}
\caption{Group differences in functional-system connectivity and the effect of motion
matching. (a, b) Standardised difference for every pair of systems in all participants and
in within-site pairs matched for head motion, on a shared colour scale; asterisks mark
pairs with Benjamini--Hochberg $q<0.05$. (c) The effect of each somatomotor region in all
participants plotted against its effect in motion-matched pairs; points below the diagonal
are stronger once motion is matched. (d) Discrimination by full-scale IQ band among
autistic participants (below 85, 85--99, 100--114, and 115 or above), each against all 468
non-autistic participants, under both validation protocols.}
\label{fig:group-differences}
\end{figure}

The gradient is reproduced when both groups are stratified by IQ, but this design is also
affected by the group difference in IQ. Specifically, all participants are divided into
quartiles of
site-centred full-scale IQ (Table~\ref{tab:iq-quartiles}). Under pooled ten-fold validation,
AUROC then falls from 0.833 in the lowest quartile to 0.709 in the highest. However,
autistic participants have lower full-scale IQ (Table~\ref{tab:participants}), so the lower
quartiles are also those in which the groups differ most on IQ. For this reason, each
autistic participant is matched to a non-autistic participant at the same site whose
full-scale IQ differs by at most five points. As a result, 274 pairs are obtained under
pooled ten-fold validation and 275 pairs with the site withheld.

\begin{table}[t]
\centering
\caption{AUROC by quartile of site-centred full-scale IQ, in all participants and in
participants matched on full-scale IQ within site. The matched counts are given for pooled
ten-fold and leave-one-site-out (LOSO) validation, respectively.}
\label{tab:iq-quartiles}
\small
\begin{tabular}{@{}lcccccc@{}}
\toprule
 & \multicolumn{3}{c}{All participants} & \multicolumn{3}{c}{IQ-matched participants} \\
\cmidrule(lr){2-4}\cmidrule(l){5-7}
Quartile & $n$ & Pooled ten-fold & LOSO & $n$ (pooled / LOSO) & Pooled ten-fold & LOSO \\
\midrule
Q1 & 204 & 0.833 & 0.825 & 137 / 144 & 0.855 & 0.848 \\
Q2 & 203 & 0.794 & 0.793 & 137 / 132 & 0.700 & 0.735 \\
Q3 & 201 & 0.746 & 0.745 & 138 / 137 & 0.751 & 0.737 \\
Q4 & 203 & 0.709 & 0.705 & 136 / 137 & 0.756 & 0.756 \\
\bottomrule
\end{tabular}
\end{table}

After IQ matching, discrimination remains highest in the lowest quartile but no longer
declines across the upper three. Specifically, AUROC in the lowest quartile is 0.855 under
pooled ten-fold validation and 0.848 with the site withheld. In contrast, AUROC across the
upper three quartiles ranges from 0.700 to 0.756 under pooled ten-fold validation and from
0.735 to 0.756 with the site withheld. Therefore, part of the gradient over the full IQ
range reflects the group difference in IQ. The within-group association, which involves no
comparison between groups, \reviewhl{is examined next as complementary evidence}.

Within the autistic group, classifier scores decline with cognitive ability, whereas no
such association is found in non-autistic participants (Table~\ref{tab:iq-correlations}). For this
analysis, site, age and mean framewise displacement are regressed from both the score and
the IQ measure (Section~\ref{subsubsec:correlation-analyses}). Under this adjustment, the score correlates with
full-scale IQ at $\rho=-0.180$ ($p=0.0004$, $n=376$) among autistic participants under
pooled ten-fold validation. In contrast, the corresponding value among non-autistic
participants is $+0.048$ ($p=0.3208$, $n=435$). In addition, the same pattern is obtained
for verbal and performance IQ and under both validation protocols.

The interaction between diagnosis and cognitive ability is also tested directly, so that the
difference between the two within-group associations is assessed in one model rather than
inferred from their separate significance. In a model of the out-of-fold score containing IQ,
diagnosis, their interaction, age, sex, mean framewise displacement and site, the interaction
is negative under both protocols. Across the six-test IQ family, it survives
Benjamini--Hochberg correction for full-scale IQ ($\beta=-0.176$, $q=0.025$) and verbal IQ
($\beta=-0.205$, $q=0.025$) under pooled ten-fold validation, and for both measures with a
site withheld ($q=0.033$). For performance IQ, however, the interaction does not survive
correction under either protocol ($q=0.082$). The dependence of the score on cognitive ability
is therefore specific to the autistic group for full-scale and verbal IQ, whereas it is not
established for performance IQ.

\begin{table}[t]
\centering
\caption{Spearman correlation between classifier score and IQ within each diagnostic group,
given as $\rho$ ($p$). Both variables are residualised on site, age and mean framewise
displacement before correlation. The p-values are not adjusted for multiplicity.}
\label{tab:iq-correlations}
\small
\setlength{\tabcolsep}{4pt}
\begin{tabular}{@{}lccccc@{}}
\toprule
 & & \multicolumn{2}{c}{Autistic} & \multicolumn{2}{c}{Non-autistic} \\
\cmidrule(lr){3-4}\cmidrule(l){5-6}
Measure & $n$ (autistic / non-autistic) & Pooled ten-fold & LOSO & Pooled ten-fold & LOSO \\
\midrule
Full-scale IQ & 376 / 435 & $-0.180$ (0.0004) & $-0.179$ (0.0005) & $+0.048$ (0.3208) & $+0.022$ (0.6466) \\
Verbal IQ & 342 / 379 & $-0.182$ (0.0007) & $-0.189$ (0.0004) & $+0.039$ (0.4547) & $+0.001$ (0.9892) \\
Performance IQ & 344 / 391 & $-0.152$ (0.0047) & $-0.147$ (0.0062) & $+0.007$ (0.8862) & $-0.005$ (0.9271) \\
\bottomrule
\end{tabular}
\end{table}

The absence of an association in non-autistic participants constrains the interpretation of
this result. Because no gradient is found in that group, the association is not attributed
to a general relation between ability and classifier score. However, the two correlations
are not compared by a direct test, and full-scale IQ varies less among non-autistic
participants (Table~\ref{tab:participants}). For this reason, this interpretation remains
provisional. In addition, \citet{Nielsenetal2013} report that classifier scores covary with
verbal IQ in ABIDE data, so the association is anticipated. Beyond that report, the
comparison with non-autistic participants and the graded loss of detectability across IQ
bands are added by the present analysis.

Taken together, these results indicate that measured performance on ABIDE-I depends on the
cognitive composition of the autistic sample. \reviewhl{The exploratory findings do not establish an IQ-specific effect because the within-group correlations are not directly compared and the IQ-matched upper quartiles are non-monotonic.} This dependence matters because
a cohort with a different distribution of cognitive ability \reviewhl{yields} a different AUROC
from the same method.

No deficit in discrimination is found in female participants, although this subgroup is
small.
Specifically, AUROC is 0.827 (0.760--0.895, $n=144$) in female participants and 0.764
(0.730--0.798, $n=727$) in male participants under pooled ten-fold validation. Similarly,
with the site withheld, AUROC is 0.806 (0.735--0.878) in female participants and 0.767
(0.733--0.801) in male participants. However, the intervals overlap, so no difference
between the sexes is claimed.

\subsection{Functional-system differences under motion matching}
\label{subsec:results-system-differences}

Group differences are assessed in exploratory analyses (Sections~\ref{subsec:group-differences}
and~\ref{subsubsec:analysis-status}), and the strengthening of the somatomotor reduction under motion
matching is the second contribution of the study. These differences are compared at the
level of functional systems, because the 19,110 individual connections are too many to
test one at a time with useful power. For this purpose, the regions are grouped into the seven cortical systems of \citet{Yeoetal2011}
and a subcortical and cerebellar class (Section~\ref{subsubsec:atlas-signal-quality}). As a result, the
comparisons are reduced to 36 system pairs, each averaging over 21 to 1,462 connections
(median 441). In addition, every connection is first adjusted for acquisition site, head
motion, age and sex. Significance is then obtained by permuting diagnosis within each site
(Sections~\ref{subsubsec:connection-adjustment}, \ref{subsubsec:system-contrasts} and~\ref{subsubsec:permutation-analyses}).

In the 863 participants of the group-contrast cohort, 10 of 36 system pairs differ after
correction, and every one of them is lower in autism (Table~\ref{tab:system-pairs};
Figure~\ref{fig:group-differences}a). Moreover, 34 of 36 pairs are negative overall, and no pair
is significantly raised. Among these pairs, the largest difference is found within the
somatomotor system (standardised difference $-0.215$, $z=-4.769$, $q=0.002$). In addition,
the limbic system is the most widely involved and differs from most other systems. Because
site enters both the adjustment and the permutation, differences in mean connectivity
between acquisition sites are removed before this reduction is estimated.

\begin{table}[t]
\centering
\caption{Functional-system pairs with the largest standardised group differences in the
group-contrast cohort (863 participants). Negative values indicate lower connectivity in
autistic participants, and $q$ is the Benjamini--Hochberg-adjusted permutation p-value.
The first ten rows are the pairs with $q<0.05$.}
\label{tab:system-pairs}
\small
\begin{tabular}{@{}llcccc@{}}
\toprule
System & System & Connections & Difference & $z$ & $q$ \\
\midrule
Somatomotor & Somatomotor & 210 & $-0.215$ & $-4.769$ & 0.002 \\
Limbic & Visual & 210 & $-0.140$ & $-3.633$ & 0.002 \\
Limbic & Salience / ventral attention & 147 & $-0.133$ & $-3.247$ & 0.013 \\
Default mode & Default mode & 903 & $-0.127$ & $-3.777$ & 0.002 \\
Default mode & Limbic & 301 & $-0.125$ & $-3.677$ & 0.004 \\
Limbic & Limbic & 21 & $-0.118$ & $-3.052$ & 0.014 \\
Dorsal attention & Somatomotor & 399 & $-0.116$ & $-2.653$ & 0.024 \\
Dorsal attention & Dorsal attention & 171 & $-0.115$ & $-2.875$ & 0.024 \\
Limbic & Somatomotor & 147 & $-0.114$ & $-2.839$ & 0.024 \\
Frontoparietal control & Limbic & 147 & $-0.102$ & $-2.733$ & 0.024 \\
Somatomotor & Visual & 630 & $-0.101$ & $-2.317$ & 0.058 \\
Frontoparietal control & Salience / ventral attention & 441 & $-0.093$ & $-2.317$ & 0.058 \\
\bottomrule
\end{tabular}
\end{table}

The permutation certifies the p-value on residualised connections, but it does not certify
the residualisation itself. For this reason, within-site pairs matched for mean framewise
displacement are also analysed (Section~\ref{subsubsec:motion-matching}). This matching is needed because
head motion alters measured connectivity \citep{Poweretal2012}, and autistic participants
move more than non-autistic participants (Table~\ref{tab:participants}). Under the
system-level caliper, 256 pairs are retained, and the group difference in motion is removed
($p=7.6\times10^{-7}$ before matching and $p=0.59$ after). In these matched pairs, 26 of 36
system pairs are significant, against 10 of 36 in the full group-contrast cohort
(Figure~\ref{fig:group-differences}b). However, the matched sample also differs from the full
cohort in size and composition, so the rise is not attributed to motion matching alone. In
addition, the somatomotor within-system difference strengthens from $-0.2155$ ($z=-4.77$)
in all participants to $-0.2925$ ($z=-4.89$) in the matched pairs. Age, however, is not
matched and is handled by adjustment alone, so the matched analysis addresses head motion
rather than age.

Inside the somatomotor system, the effect is spread across its 21 regions rather than
concentrated in one (Supplementary Material). Specifically, all 21 regions differ
significantly in all participants, and the strongest effects lie in the postcentral gyrus,
the paracentral lobule and the precentral gyrus. In anatomical terms, these regions form
primary somatosensory and motor cortex, and they are affected on both sides to a similar
degree. However, this regional analysis is exploratory, because the system is chosen for
carrying the largest system-level difference (Section~\ref{subsubsec:regional-analysis}). Its ordering is,
therefore, a description of that result rather than an independent test of it.

The regional result also \reviewhl{persists after matching on recorded mean FD and is larger in the matched sample}.
Specifically, 280 pairs (560 participants) are retained by the wider regional caliper, and
the motion difference is again removed ($p=0.28$ after matching). Under this matching, all 21
regions remain significant, every one of them strengthens, and their ordering is preserved
(rank correlation 0.89 against the unmatched result). Moreover, the mean regional effect moves from $-0.215$ in all 863
participants to $-0.249$ in the matched pairs
(Figure~\ref{fig:group-differences}c).

The same conclusions are obtained when the tests are repeated under the Freedman--Lane
residual permutation, which holds the association between diagnosis and the covariates fixed
(Section~\ref{subsubsec:permutation-analyses}). Under that procedure, the number of system
pairs with reduced connectivity and $q<0.05$ rises from 14 to 27 of 36 after motion matching,
and every one of those pairs is negative. The somatomotor pair itself strengthens from
$-0.226$ to $-0.294$. In addition, all 21 somatomotor regions remain significant at $q<0.05$
in both the full and the matched analyses, and the mean regional effect strengthens from
$-0.226$ to $-0.294$. The direction, the strengthening under motion matching, and the
involvement of every region are therefore not artefacts of the exchangeability assumption
used in the primary tests.

The regional ordering is moderately reproducible within the cohort, and this reproducibility
\reviewhl{indicates the reliability of the regional ranking}. Specifically, a correlation of
0.55 across the 21 regions is obtained in repeated split-half analyses, with the nuisance
adjustment refitted in each half. In comparison, the ordering of the 210 connections inside the
system reproduces at 0.38, and that of its top fifty at 0.34. In addition, 100 random
blocks of 21 regions drawn from the whole brain give a median of 0.35 (range 0.06--0.63),
and 9\% of them reach the somatomotor value. The ordering is, therefore, more stable than
most blocks of equal size, but not categorically so.

The reduction runs against the direction expected from residual motion. Specifically, mean
framewise displacement correlates positively with mean connection strength ($\rho=+0.51$
across participants and $+0.45$ within site). Higher motion is, therefore, associated with
higher rather than lower connectivity in this pipeline. Because autistic participants move
more, residual motion would on average be expected to raise their connectivity, which is
opposite to the observed reduction. In addition, the somatomotor location of the difference
agrees with the somatomotor involvement reported by \citet{Chenetal2015}. Its direction
also agrees with the somatomotor reductions reported by \citet{Holigaetal2019} and
\citet{Ilioskaetal2023} in samples that include ABIDE data. Because those samples \reviewhl{could}
share participants with the present cohort, this agreement is not an independent
replication. The present result, in contrast, is obtained in within-site pairs matched for
head motion.

These group differences carry three limits. First, residual motion is not a sufficient
explanation for them, but it is not thereby excluded as a contributor. Second, the
contrasts are univariate and are computed without reference to the classifier, so they
support no statement about which measurements the classifier uses. Finally, all group
differences are obtained under a single denoising strategy without global signal
regression, so their direction is conditional on that choice.

\subsection{Feature-selection scope and granularity}
\label{subsec:results-selection-scope}

The effect of data-access scope on measured performance is examined in seven conditions.
Among these conditions, the inflation produced by feature selection outside the training
partition, and its dependence on feature granularity, form the third contribution. For
this purpose, the conditions are compared on identical folds with the 19,900 static tangent
features of the primary cohort (Section~\ref{par:selection-scope}; Table~\ref{tab:selection-scope}). Under
this design, fold-mean AUROC is changed by at most 0.003 when the tangent reference, or the
reference and the scaler, is fitted on all participants. Similarly, it is changed by a
similar amount when the regularisation constant is chosen on the held-out fold (0.758) or
when acquisition site is added as a predictor (0.756).

\begin{table}[t]
\centering
\caption{Fold-mean AUROC under the seven data-access conditions, with 19,900 static tangent
features in the primary cohort (871 participants) and identical outer folds. SD is the
standard deviation across outer folds.}
\label{tab:selection-scope}
\small
\begin{tabular}{@{}>{\raggedright\arraybackslash}p{0.62\linewidth}cc@{}}
\toprule
Condition & AUROC & SD \\
\midrule
Reference: every step fitted within the training partition & 0.756 & 0.048 \\
Tangent reference fitted on all participants & 0.756 & 0.048 \\
Tangent reference and scaler fitted on all participants & 0.753 & 0.049 \\
Regularisation constant chosen on the held-out fold & 0.758 & 0.048 \\
Acquisition site added as a predictor & 0.756 & 0.049 \\
Tangent reference fitted on all participants and 2,000 of 19,900 connections selected with all diagnostic labels & 0.987 & 0.008 \\
All five changes combined & 0.988 & 0.008 \\
\bottomrule
\end{tabular}
\end{table}

In contrast, fold-mean AUROC is raised from 0.756 to 0.987 when 2,000 of 19,900 connections
are selected with an F test on all diagnostic labels before splitting. Likewise, accuracy
rises from 0.692 to 0.948 on the same folds. Because fitting the tangent reference on all
participants leaves AUROC at 0.756 on its own, this rise is attributed to the selection
step. As a result, an inflation of 0.231 in AUROC is produced by this step, and this
inflation is the largest effect in the study. Moreover, 0.988 is obtained when every change
is applied together, so the remaining changes add little beyond the selection step.
Therefore, of the changes examined, only feature selection with all labels changes the
score by more than 0.003. By comparison, choosing the regularisation constant on the
held-out fold also uses held-out labels, but its effect lies within 0.003. As a
consequence, figures from pipelines that select connections before cross-validation are
not comparable with estimates in which selection is nested.

Selection at the level of whole regions, in contrast, gives no consistent advantage
(Table~\ref{tab:region-selection}; Figure~\ref{fig:corrections}c). For this comparison, regions
are ranked either with all diagnostic labels or with the labels of each training fold only
(Section~\ref{par:selection-scope}). Under this design, the two scopes differ by 0.006 on average in
absolute value across 10 to 100 of 196 regions. Specifically, the difference is 0.001 at 10
regions and at most 0.012 at 100 regions, and its sign is inconsistent. Therefore, the risk
of a selection step depends on feature granularity. In particular, a choice of 10 of 196
regions is too coarse to fit noise, whereas a choice of 2,000 of 19,900 connections is not.

This conclusion does not depend on the scope of the nuisance adjustment used for the ranking.
When the nuisance model is refitted within each outer-training partition, so that no held-out
participant contributes to the residuals from which the regions are ranked
(Section~\ref{par:selection-scope}), the AUROC of the training-fold condition changes by 0.0035
on average in absolute value, and by at most 0.0085 at any list length. The ordering across
list lengths is preserved, and every restricted model remains below the model that keeps all
196 regions (0.781). The comparison is therefore unchanged by making the selection strictly
nested.
However, the two comparisons also differ in cohort (871 against 863 participants),
additional shrinkage, ranking statistic and AUROC summary (Section~\ref{par:selection-scope}). For
this reason, the contrast is consistent with an effect of granularity, but granularity is
not the only difference between the two comparisons.

\begin{table}[t]
\centering
\caption{Participant-pooled AUROC when the top-ranked regions are chosen with all
diagnostic labels or with the labels of each training fold only, in the group-contrast
cohort (863 participants and 196 regions). All 196 regions give 0.781.}
\label{tab:region-selection}
\small
\begin{tabular}{@{}ccc@{}}
\toprule
Regions kept & Chosen with all labels & Chosen within each fold \\
\midrule
10 & 0.628 & 0.627 \\
20 & 0.660 & 0.656 \\
30 & 0.699 & 0.702 \\
50 & 0.738 & 0.748 \\
100 & 0.764 & 0.752 \\
\bottomrule
\end{tabular}
\end{table}

\begin{figure}[pos=tbp]
\centering
\includegraphics[width=\linewidth]{fig_corrections}
\caption{Four effects that change when the corresponding control is applied. (a)
Correlation between classifier score and symptom severity under pooled ten-fold validation,
computed across both groups and within the autistic group; rows marked with a dagger are
instruments given only to autistic participants, for which the two scopes necessarily coincide.
(b) AUROC of whole-scan connectivity across shrinkage levels, with the windowed
representation placed at the shrinkage that it applies to its own windows. (c) AUROC
inflation from selection outside the training partition, for 10 to 100 of 196 regions and
for 2,000 of 19,900 connections. (d) Observed accuracy against the interval of accuracy
attainable given the diagnostic labels.}
\label{fig:corrections}
\end{figure}

Restricting the model to the most discriminating regions also lowers performance at every
list length. Specifically, all 196 regions give an AUROC of 0.781, which is higher than any
reduced set in Table~\ref{tab:region-selection}. In particular, the smallest set loses 0.154
against the use of every region. This pattern, in turn, suggests that the information used
for discrimination is spread across many regions rather than confined to those with the
largest group differences. Consequently, a reduced set of regions does not improve
discrimination on these data, and this restriction is reported as a limit.

Without any selection before the split, the published accuracy range is already reached by
the ensemble. Specifically, a participant-pooled accuracy of 0.706 is obtained in the wider
cohort under pooled ten-fold validation (Section~\ref{subsec:results-performance}). This value is
comparable to the 0.700, 0.703 and 0.708 reported for the same eligible count of 1,035
participants. Accuracy at that level, therefore, does not require selection
outside the training partition.

\subsection{Classifier family}
\label{subsec:results-classifier-family}

Linear classifiers outperform the fixed-configuration tree ensembles on static tangent
connectivity, and this result is the fourth contribution. For this comparison, all eight
classifiers use the same features, feature scaling and outer folds
(Section~\ref{par:classifier-family}), so differences are not confounded by the input or the partition.
Under this design, participant-pooled AUROCs of 0.735--0.757 are reached by the five linear
classifiers, whereas 0.611--0.624 is reached by random forest, XGBoost and CatBoost
(Table~\ref{tab:classifier-family}). In particular, 0.757 is reached by penalised logistic
regression, and ridge (0.755) and linear support-vector classification (0.753) lie close to
it.
In contrast, shrinkage linear discriminant analysis (0.741) and the rank-20 spectral model
(0.735) are lower.

\begin{table}[t]
\centering
\caption{Classifier family on the static tangent features of the primary cohort
(871 participants), with identical features, scaling and outer folds. Only logistic
regression is tuned; the other classifiers use fixed configurations.}
\label{tab:classifier-family}
\small
\setlength{\tabcolsep}{4pt}
\begin{tabular}{@{}lcccc@{}}
\toprule
Classifier & AUROC (participant-pooled) & AUROC (fold mean) & SD (folds) & Accuracy (fold mean) \\
\midrule
Logistic regression & 0.757 & 0.756 & 0.048 & 0.692 \\
Ridge & 0.755 & 0.755 & 0.047 & 0.693 \\
Linear support-vector & 0.753 & 0.753 & 0.048 & 0.687 \\
Shrinkage LDA (PCA-200) & 0.741 & 0.739 & 0.052 & 0.669 \\
Low-rank spectral (rank 20) & 0.735 & 0.737 & 0.044 & 0.672 \\
CatBoost & 0.624 & 0.623 & 0.042 & 0.597 \\
XGBoost & 0.617 & 0.616 & 0.050 & 0.584 \\
Random forest & 0.611 & 0.609 & 0.049 & 0.577 \\
\bottomrule
\end{tabular}
\end{table}

The gap between logistic regression and each tree ensemble is large relative to its
uncertainty. Specifically, the participant-pooled AUROC of logistic regression is higher
than that of XGBoost by 0.1396, of random forest by 0.1461 and of CatBoost by 0.1327. For these differences,
unadjusted DeLong p-values range from $6.3\times10^{-15}$ to $2.4\times10^{-12}$
(Section~\ref{subsubsec:paired-comparisons}). By comparison, the rank-20 spectral model lies 0.0214 below
logistic regression (unadjusted $p=0.0439$; $q=0.1065$, for reference only).

This contrast characterises the configurations evaluated rather than the best attainable
member of each family. Specifically, only the logistic model is tuned, whereas the tree
ensembles use fixed configurations without nested tuning (Section~\ref{par:classifier-family}).
However, the fixed-configuration ridge, linear support-vector and discriminant classifiers
also exceed all three tree ensembles. Therefore, the advantage of these linear
configurations over these tree-ensemble configurations does not rest on the tuning of the
logistic model alone.

\subsection{Scope of the symptom-severity associations}
\label{subsec:results-symptom-severity}

Under pooled ten-fold validation, no association between the classifier score and symptom
severity survives correction within the autistic group, and this absence forms the
fifth contribution of the study (Figure~\ref{fig:corrections}a). For this question, the correlation is computed in
two scopes, namely across both groups and within the autistic group only. The second scope
is needed because non-autistic participants who complete these instruments score at the
low end of the scale. In this cohort, for example, 19 non-autistic participants have an
ADOS total and 137 an SRS total. For this reason, any model that separates the groups
produces a correlation across both groups largely by construction. Consistent with this, 3
of 6 associations survive Benjamini--Hochberg correction across both groups, whereas 0 of 6
do within the autistic group (Section~\ref{par:symptom-associations}).

The largest association across both groups almost vanishes once non-autistic participants
are removed. Specifically, for the SRS total, a correlation with the classifier score of
$\rho=0.438$ is obtained across both groups ($n=300$), but of 0.050 within the autistic
group ($n=163$). Similarly, values of 0.207 against 0.115 are obtained for the ADOS total,
and of 0.196 against 0.128 for the ADOS calibrated severity. In contrast, identical
correlations in both scopes are obtained for the ADI-R domains, which are given only to
autistic participants, because both scopes contain the same people. Therefore, the
correlations across both groups reflect the separation of the groups. An association with
the degree of autism within the autistic group, in turn, is not established by these data.

\subsection{Windowed connectivity and covariance shrinkage}
\label{subsec:results-windowed-shrinkage}

The small advantage of the windowed representation is not distinguished from chance, and it
is not retained when the whole-scan estimate receives the same covariance shrinkage; this
result is the sixth contribution (Figure~\ref{fig:corrections}b). Specifically, an extra shrinkage of 0.15 is applied by
the windowed representation to its window covariances, and this shrinkage is not received
by the whole-scan representation (Section~\ref{subsubsec:primary-comparisons}). For this reason, the same
shrinkage is applied to the whole-scan estimate, with nothing else changed. Under this
control, the static AUROC is raised by 0.013 through shrinkage alone (unadjusted DeLong
$p=0.004$). In contrast, the windowed representation differs from the unshrunk static
baseline by 0.010 ($p=0.281$) and from the matched static baseline by $-0.002$
($p=0.709$). Therefore, no evidence is obtained that windowing adds information beyond the
shrinkage that it applies, and the apparent benefit of windowing is attributed to that
shrinkage.

\subsection{Secondary results}
\label{subsec:results-secondary}

The remaining analyses are reported as descriptive results, and none of them is claimed as
a contribution. In the first of these, the ensemble probabilities have \reviewhl{low ten-bin expected calibration error}
under both protocols (Figure~\ref{fig:performance}c). Specifically, the
expected calibration error of Equation~\eqref{eq:calibration-error} is 0.046 under pooled ten-fold
validation and 0.040 under leave-one-site-out validation. In addition, these values are
computed from participant-pooled predictions, for which accuracy is 0.707 and 0.711,
respectively. Because this error is the bin-weighted mean gap between predicted probability
and observed frequency, small values indicate agreement on average. However, a small error
does not guarantee agreement within every probability bin \reviewhl{or establish calibration without slope, intercept, Brier-score, and uncertainty analyses}, so the reliability diagram is
also shown.

Measures that do not depend on binning give the same picture, with one exception. The
calibration slope is 1.117 (95\% site-clustered bootstrap interval 0.866 to 1.321) under
pooled ten-fold validation and 1.119 (0.856 to 1.349) with a site withheld, so both intervals
include one and the probabilities are neither over- nor under-dispersed. The Brier score is
0.194 under both protocols. Calibration-in-the-large, however, is $-0.162$ ($-0.282$ to
$-0.061$) under pooled ten-fold validation, and this interval excludes zero, so the average
predicted risk is slightly above the observed rate. With a site withheld the corresponding
value is $-0.181$ ($-0.385$ to $0.019$). The binned error itself depends on the binning,
rising from 0.037 at five bins to 0.046 at ten and 0.059 at twenty, where some bins contain no
participants. The expected calibration error is therefore reported with its bin count rather
than as a single figure.

Calibration, however, does not translate into a useful guarantee on the cases that the
model chooses to answer (Figure~\ref{fig:performance}d; Figure~\ref{fig:abstention-guarantee}). To
test such a guarantee, the leave-one-site-out predictions are divided at random into two
halves, and the threshold is chosen on one half and applied to the other
(Section~\ref{subsubsec:threshold-analysis}). This division is needed because a threshold scored on the
participants that set it passes its bound by construction. When the cohort-wide threshold
is set on the whole sample, however, no participant is answered. In addition, under the
half-split test, the target is met in a proportion of 0.8158 of the 38 half-splits with a
finite selective risk, which is below the nominal 0.90.

\begin{figure}[pos=tbp]
\centering
\includegraphics[width=\linewidth]{fig_guarantee}
\caption{The error guarantee of the abstention rules under leave-one-site-out validation and
its evaluation on held-out participants. (a) The error bound that each rule reports on the participants used to set
it, against the error that the same rules incur on held-out cases. (b) The cost of each
rule for obtaining that bound: the proportion of cases it answers, and the proportion of
held-out splits in which it meets its target.}
\label{fig:abstention-guarantee}
\end{figure}

The target is held by the per-site threshold, but few participants are answered by it.
Specifically, when it is set on the whole sample, 170 of 871 participants (19.5\%) are
answered, and 143 of them are at a single site. Moreover, no participant is answered at 17
of the 20 sites. Across held-out half-splits, in addition, its mean coverage, the
proportion of participants answered, is 13.7\%, and the target is met in a proportion of
0.9137 of 197 valid splits. Abstention is, therefore, reported as a
measurement on these data rather than as a feature of the model.

Describing connectivity with more than one atlas does not improve discrimination either.
Specifically, a fold-mean AUROC of 0.756 is reached by CC200 alone, against 0.749 for CC200
with the anatomical labelling atlas, 0.755 for CC200 with the Harvard--Oxford atlas, and
0.752 for all three together (Section~\ref{par:atlas-fusion}). In addition, each anatomical
atlas reaches less on its own, namely 0.692 and 0.717. The functional parcellation
therefore carries the information that the combinations provide, and no gain is obtained
from adding atlases.

The diagnostic labels themselves limit attainable accuracy, but the size of this limit
cannot be fixed from these data (Figure~\ref{fig:corrections}d). To estimate it, the two
diagnostic instruments are compared in the 213 autistic participants with both a complete
ADI-R algorithm and an ADOS total (Section~\ref{par:instrument-concordance}). Under this comparison, the
instruments disagree in 15.0\% of these participants. Specifically, 13 meet criteria on the
interview only, 19 on the observation only and 181 on both, whereas no participant is
negative on both instruments. Under the
assumptions of Section~\ref{par:instrument-concordance}, the accuracy attainable against these labels
therefore lies between approximately 0.85 and 0.99. Observed accuracy, in turn, lies 0.14
below even the lower end of that range.

At the central noise rates, correction for label noise does not close this gap. These
central rates are the values defined in Section~\ref{par:instrument-concordance}, namely half the observed
disagreement and the proportion of non-autistic participants with an SRS raw total of at
least 60. Specifically, when the correction of \citet{Menonetal2015} is applied at these
rates, accuracy is moved from 0.707 to 0.729, and AUROC from 0.776 to 0.806. The shortfall
is, therefore, not explained by uncertainty in the diagnosis alone. However, at the most pessimistic rates the correction is larger, so the size of the
remaining gap depends on assumptions that these data cannot settle. In particular,
agreement beyond chance is slightly negative (Cohen's $\kappa=-0.078$), which is
inconsistent with the two-rater model that the correction assumes. Moreover, two
instruments applied to one population do not identify their error rates, and, as noted above, no participant is negative on both. For these
reasons, no point
estimate of the ceiling is reported.
